% Solutions to the exercises of Chapter 24.
\begin{solutions}
  \item \emph{EKF-SLAM, 200 landmarks:} $n = 15 + 600 = 615$, so $P$ has $378\,225$ entries, 3.0~MB in double precision. One image observation has a Jacobian with 18
        non-zero columns, so $PH\T$ costs about $2\cdot615\cdot18\cdot2 \approx 44$~kflop, but the covariance update $P - KSK\T$ touches every entry, about $4n^2 \approx 1.5$~Mflop. With 100
        observations per frame this is about 160~Mflop per frame, 3~Gflop/s at 20~Hz, and it grows with the square of the map. (A dense Joseph form, $\approx 4n^3$, would cost
        almost 1~Gflop per update.) \emph{MSCKF, 10 clones:} $n = 75$, so $P$ has 5\,625 entries, 45~kB, small enough to stay in cache. An update with 50 features of 6
        observations stacks $50\cdot(12 - 3) = 450$ rows after nullspace projection (about $0.3$~Mflop for the projections). QR compression to 75 rows costs about $4.8$~Mflop,
        and the EKF update with 75 rows, Joseph form included, about 6~Mflop: roughly 11~Mflop per update. The memory is 70 times smaller and the flops more
        than ten times fewer. More importantly, the MSCKF's cost does not depend on how many features are being tracked, only on how many are processed.

  \item The $2M$ scalar observations involve the 3 unknown coordinates of the feature, so after eliminating the feature $2M - 3$ constraints remain on the poses.
        With $M = 1$ there are none: the two numbers do not even determine the feature, whose depth stays free. With $M = 2$ there is one, the epipolar
        (coplanarity) constraint between the two poses (Exercise~26.3). Each further view adds two. The project requires $M \ge 3$ because a two-view track is a weak and
        fragile measurement. Its triangulation rests on the baseline between two consecutive images, often a few centimeters, so the linearization point
        is poor. Its single projected residual gives the $\chi^2$ gate one degree of freedom, too little to reject a bad match. And the track has no redundancy
        with which to check the triangulation itself. The information it would contribute is small, and the risk is not.

  \item The past estimate $\hat x_1$ and the current estimate $\hat x_2$ were produced by the same propagation from the same measurements, so their errors are correlated. Treating
        $\hat x_1$ as exact turns a \emph{relative} measurement into an \emph{absolute} one, credited with an accuracy the anchor does not have. Scalar example: $x_1 \sim \N(\hat x_1, \sigma_1^2)$,
        $x_2 = x_1 + w$ with $w \sim \N(0, q)$, and $z = x_2 - x_1 + v$ with $v \sim \N(0, r)$; take $\sigma_1^2 = q = r = 1$.
        \begin{itemize}
          \item \emph{Correct (cloning):} $x_1$ is kept as a state, $H = [-1\ \ 1]$, and the posterior covariance of $(x_1, x_2)$ is $\big[\begin{smallmatrix} 1 & 1 \\ 1 & 1.5\end{smallmatrix}\big]$.
                The variance of $x_2$ drops from 2 to 1.5, and that of $x_1$ is unchanged.
          \item \emph{Wrong:} $z + \hat x_1$ is used as a measurement of $x_2$ with variance $r$. The filter claims $\Var(x_2) = 2\cdot1/3 = 0.67$, but the actual error
                variance is $\sigma_1^2 + (1 - K)^2q + K^2r = 1.56$ with $K = \tfrac23$ (a Monte Carlo agrees): overconfident by a factor 2.3, and less accurate than the correct filter.
        \end{itemize}
        Repeating such anchored updates makes it worse: after 20 steps the claimed variance has settled at $0.62$ while the actual error variance has grown to
        11.6, like a random walk. The filter believes it is localized when it is drifting.
\end{solutions}
